% !TeX spellcheck = pt_BR
\documentclass[a4paper,12pt]{letter}  % tipo do documento
%\usepackage[brazil]{babel}           % define a linguagem padrão
\usepackage{graphicx}                % permite a inserção de figuras
\usepackage[T1]{fontenc}
% \usepackage[latin1]{inputenc}        % para usar caracteres acentuados diretamente
\usepackage[utf8]{inputenc}
\usepackage[hmargin=1.5cm,vmargin=1cm]{geometry}
\usepackage{listings} %para colocar codigos
\usepackage{amsmath}
\usepackage{amssymb}
\usepackage{multicol}
\setlength{\columnsep}{1.5cm}
\setlength{\columnseprule}{0.2pt}

% \usepackage{splitidx}
% \lstset{numbers=left,
% stepnumber=1,
% firstnumber=1,
% numberstyle=\tiny,
% extendedchars=true,
% breaklines=true,
% frame=tb,
% basicstyle=\footnotesize,
% % stringstyle=\ttfamily,
% % showstringspaces=false
% }




% \newcommand{\gabarito}[1]{\textbf{\begin{small} \\ #1 \\ \end{small}}} %mostrar gabarito
\newcommand{\gabarito}[1]{\textbf{\begin{small} \end{small}}} %nao mostrar gabarito


\begin{document}
\thispagestyle{empty}
% \pagestyle{empty}

UNIVERSIDADE FEDERAL DO CEARÁ\\
Departamento de Estatística e Matemática Aplicada\\
Disciplina de Elementos de Análise Combinatória cc0279.\\
Professor: Albert Einstein F. Muritiba.\\
Temas: Lógica e Conjuntos\\[0.2cm]
\textbf{Nome}: \makebox[380pt]{\hrulefill} \vspace{2pt}
%Matr\'icula:\makebox[60pt]{\hrulefill} \vspace{2pt}
\begin{center}
	1º AP
	% -- \textit{2ª Chamada}
	% 2\textordfeminine Avaliação - Segunda Chamada
	% npc 2 %+ nef
	% NEF
\end{center}
%\begin{multicols}{2}
Responda [V] para Verdadeiro e [F] para Falso.
\begin{enumerate}
    \item Da afirmação "Se Maria for à praia, então ela irá ao cinema", podemos concluir que:
    \begin{enumerate}
        \item ( ) Maria vai à praia ou não vai ao cinema.
        \item ( ) Se Maria for ao cinema, então ela irá à praia.
        \item ( ) Se Maria não for à praia, então ela não irá ao cinema.
        \item ( ) Se Maria não for ao cinema, então ela não irá à praia.
        \item ( ) Maria ir à praia é condição necessária e suficiente para ela ir ao cinema.
    \end{enumerate}

    \item Sejam $A$ e $B$ conjuntos quaisquer e $A \subset B$. Podemos \textbf{sempre} afirmar que:
    \begin{enumerate}
        \item ( ) $(\forall x)[x \in A \land x \in B]$
        \item ( ) $(\forall x)[x \in A \rightarrow x \in B]$
        \item ( ) $(\exists x)[x \in A \lor x \in B]$
        \item ( ) $(\forall x)[x \notin A \rightarrow x \notin B]$
        \item ( ) $(\forall x)[x \notin B \rightarrow x \notin A]$
    \end{enumerate}

    \item A afirmação "Não é verdade que todos os alunos do curso de Estatística são inteligentes ou ricos" é equivalente a quais das seguintes afirmações? ( V-equivalente, F-não equivalente)
    \begin{enumerate}
        \item ( ) "Todos os alunos do curso de Estatística não são inteligentes e ricos."
        \item ( ) "Existe aluno do curso de Estatística que não é inteligente ou é rico."
        \item ( ) "Algum aluno do curso de Estatística não é inteligente ou não é rico."
        \item ( ) "Há aluno do curso de Estatística que não é inteligente e não é rico."
        \item ( ) "Todos os alunos do curso de Estatística não são inteligentes ou não são ricos."
    \end{enumerate}

    \item Sejam os conjuntos $X=\{0,1,2\}$, $Y = \{0,1,2,\{0,1,2\}\}$, $Z=\{\emptyset, X\}$ . 		      
        \begin{enumerate}
            \item ( ) $Z \cap Y = X$
            \item ( ) $X \in Y \wedge X \subsetneq Y$
            \item ( ) $X - Z = \emptyset$
            \item ( ) $\#Y - \#Z = \#(X-Y)$  -- ($\#A$ é o número de elementos do conjunto $A$)
            \item ( ) $ \emptyset \subset Z$
        \end{enumerate}

    \item Demonstre usando:
    \begin{enumerate}
        \item Raciocínio dedutivo: 
         $$  (p \rightarrow q) \land (p \rightarrow \sim q) \Rightarrow \sim p $$
        \item Validade de argumentos: 
            \begin{align*}
	            H_1:& &B \lor C \rightarrow  B \land A,\\
	            H_2:& &\sim B,\\
	            \therefore & & \sim C.
	        \end{align*}
    \end{enumerate}
           
\end{enumerate}
\end{document}